\documentclass[12pt,a4paper]{article}

% ============================================================
% PACOTES
% ============================================================
\usepackage[utf8]{inputenc}
\usepackage[T1]{fontenc}
\usepackage[brazil]{babel}
\usepackage{amsmath,amssymb,amsfonts,mathtools}
\usepackage{bm}
\usepackage{booktabs}
\usepackage{geometry}
\usepackage{graphicx}
\usepackage{float}
\usepackage{hyperref}
\usepackage{listings}
\newcommand{\Prob}{\mathbb{P}}
\usepackage{pgfplots}
\pgfplotsset{compat=1.18}
\usepackage{tikz}

\geometry{a4paper, margin=2.5cm}

\hypersetup{
    colorlinks=true,
    linkcolor=blue,
    citecolor=blue,
    urlcolor=blue
}

% ============================================================
% COMANDOS
% ============================================================
\newcommand{\avg}[1]{\left\langle #1\right\rangle}
\newcommand{\e}[1]{\mathrm{e}^{#1}}
\newcommand{\Var}{\mathrm{Var}}
\newcommand{\nn}{\mathrm{nn}}
\newcommand{\thetaT}{\theta}

\begin{document}
% ============================================================
% CABEÇALHO
% ============================================================
\begin{center}
    {\Large \textbf{Sexta série de exercícios}}\\[0.6cm]
    {\large Tópicos de Mecânica Estatística -- Transições de Fases -- 2026}\\[0.2cm]
    {\large Urna de Ehrenfest e Monte Carlo -- 13/10/2026}\\[0.2cm]
    {\large Data de Entrega -- 03/11/2026}\\[0.2cm]
    {\Large \textbf{Autor:} Thiago Siqueira Domingues (Nº USP: 8942194)}
\end{center}
%
\vspace{0.4cm}
% ============================================================
% PROBLEMA 1
% ============================================================
\section*{Problema 1 -- Modelo das urnas de Ehrenfest}
%
Há $N$ bolas numeradas distribuídas entre duas urnas. Denotamos por
\[
N_1(k)=n,\qquad N_2(k)=N-n,
\]
onde $k=0,1,2,\ldots$ é o número de atualizações realizadas. Em cada atualização escolhe-se uma das $N$ bolas com probabilidade uniforme e transfere-se a bola escolhida para a outra urna. Portanto, exatamente uma bola muda de urna a cada passo. Se o estado atual é $N_1(k)=n$, então a probabilidade de selecionar uma bola da urna 1 é $n/N$, e a de selecionar uma bola da urna 2 é $(N-n)/N$. Consequentemente,
%
\begin{align}
\Prob\bigl(N_1(k+1)=n-1\mid N_1(k)=n\bigr)&=\frac{n}{N}, \label{eq:down}\tag{1}\\
\Prob\bigl(N_1(k+1)=n+1\mid N_1(k)=n\bigr)&=\frac{N-n}{N}. \label{eq:up}\tag{2}
\end{align}
%
Os estados admissíveis são $n=0,1,\ldots,N$. As fórmulas já respeitam os limites, se $n=0$, a probabilidade de diminuir $n$ é zero; se $n=N$, a probabilidade de aumentar $n$ é zero. Neste problema, cada passo corresponde a um intervalo de tempo $\Delta t$. Para evitar confundir uma variável temporal contínua com o índice inteiro de atualizações, usaremos $k$ nas equações discretas. Se necessário, o tempo físico é $t=k\Delta t$.
%
\section*{(a) Simulação, rota para o equilíbrio e flutuações}
%
A simulação pode ser implementada da seguinte maneira:
%
\begin{enumerate}
  \item Fixar $N$ e a condição inicial $N_1(0)=N$ (todas as bolas começam na urna 1).
  \item Em cada passo $k$, gerar $r$ uniformemente em $[0,1)$.
  \item Se $r<N_1(k)/N$, fazer $N_1(k+1)=N_1(k)-1$; caso contrário, fazer $N_1(k+1)=N_1(k)+1$.
  \item Registrar $N_1(k)$ e repetir o processo.
\end{enumerate}
%
A desigualdade estrita ou não estrita no teste não altera as probabilidades, pois a probabilidade de obter exatamente um valor real preestabelecido de $r$ é zero. Quando $N_1$ é maior que $N/2$, há mais bolas na urna 1 do que na urna 2 e, portanto, é mais provável escolher uma bola da urna 1 e diminuir $N_1$. Quando $N_1<N/2$, ocorre o contrário. Isso produz uma tendência estatística em direção a $N_1\simeq N/2$.

É importante distinguir a média de uma trajetória individual, uma trajetória não fica parada em $N/2$. Mesmo depois da relaxação, $N_1(k)$ continua mudando e flutuando. O equilíbrio é estatístico, o número de bolas em cada urna é, em média, $N/2$, com flutuações não nulas. A distribuição estacionária e os momentos dessas flutuações serão derivados nos itens (b)--(d).

A distribuição estacionária será demonstrada no item (c) como a distribuição binomial
%
\begin{equation}
  P_n=\Prob_{\mathrm{est}}(N_1=n)=\frac{1}{2^N}\binom{N}{n},\qquad n=0,1,\ldots,N.
  \label{eq:binomial}
\end{equation}
%
Ela é uma binomial com parâmetros $N$ e $p=1/2$. Assim, como será provado explicitamente no item (c),
%
\begin{equation}
  \avg{N_1}_{\mathrm{est}}=\frac{N}{2},\qquad
  \Var(N_1)=\avg{(N_1-\avg{N_1})^2}=\frac{N}{4}.
  \label{eq:meanvar_preview}
\end{equation}
%
O desvio padrão é a raiz quadrada da variância:
%
\begin{equation}
  \boxed{\sigma_{N_1}=\sqrt{\Var(N_1)}=\frac{\sqrt{N}}{2}.}
  \label{eq:sigma}
\end{equation}
%
A flutuação relativa, definida como desvio padrão dividido pelo valor médio, vale
%
\begin{align}
  \frac{\sigma_{N_1}}{\avg{N_1}_{\mathrm{est}}}
  &=\frac{\sqrt{N}/2}{N/2}
   =\frac{\sqrt{N}}{2}\frac{2}{N}
   =\boxed{\frac{1}{\sqrt{N}}}.
  \label{eq:relativefluct}
\end{align}
%
Portanto, as flutuações absolutas crescem como $\sqrt{N}$, ao passo que as flutuações relativas diminuem como $N^{-1/2}$. Este é o conteúdo da chamada lei de flutuações $\sqrt{N}$. Para os dois valores pedidos,
%
\begin{center}
\begin{tabular}{@{}cccc@{}}
\toprule
$N$ & $\avg{N_1}_{\mathrm{est}}=N/2$ & $\sigma_{N_1}=\sqrt{N}/2$ & $\sigma_{N_1}/\avg{N_1}=1/\sqrt{N}$\\
\midrule
$20$  & $10$  & $2.236\ldots$ & $0.2236\ldots\ (22.4\%)$\\
$200$ & $100$ & $7.071\ldots$ & $0.07071\ldots\ (7.07\%)$\\
\bottomrule
\end{tabular}
\end{center}
%
Assim, para $N=20$ esperamos flutuações típicas da ordem de duas bolas ao redor de $10$; para $N=200$, flutuações típicas da ordem de sete bolas ao redor de $100$. O símbolo ``flutuações típicas'' não significa que toda observação esteja necessariamente dentro de exatamente um desvio padrão.

Como a fração $N_1/N$ é uma grandeza limitada entre zero e um, é útil representar $N_1(k)/N$ em função do tempo reescalado $k/N$. A justificativa para essa escolha vem da evolução média exata, derivada no item (d): para $N_1(0)=N$,
%
\begin{equation}
  \frac{\avg{N_1}_k}{N}
  =\frac{1}{2}\left[1+\left(1-\frac{2}{N}\right)^k\right].
  \label{eq:meanfraction}
\end{equation}
%
Para $N\gg1$, usamos
%
\begin{equation}
  \ln\left(1-\frac{2}{N}\right)\simeq -\frac{2}{N},
  \qquad
  \left(1-\frac{2}{N}\right)^k
  =\exp\left[k\ln\left(1-\frac{2}{N}\right)\right]
  \simeq \exp(-2k/N).
\end{equation}
%
Logo,
%
\begin{equation}
  \frac{\avg{N_1}_k}{N}\simeq \frac{1}{2}\left(1+\exp(-2k/N)\right).
  \label{eq:meancontinuum}
\end{equation}
%
A curva média depende aproximadamente apenas de $k/N$, o que explica por que essa variável é conveniente ao comparar $N=20$ e $N=200$. A figura mostra exemplos de trajetórias simuladas e a média teórica. Uma única realização é ruidosa; o colapso é mais claro na curva média do que em cada trajetória individual.
%
\begin{figure}[htbp]
  \centering
  \includegraphics[width=0.92\linewidth]{Figures/statistical_mechanics/fig_ehrenfest_problema1.png}
  \caption{Exemplos de trajetórias simuladas para $N=20$ e $N=200$, usando $N_1(0)=N$. A linha contínua é a média teórica discreta, e a tracejada é a aproximação para $N$ grande. A trajetória permanece flutuando após a relaxação. A realização específica depende da semente do gerador de números aleatórios.}
  \label{fig:ehrenfest}
\end{figure}
%
\section*{(b) Equação de evolução temporal e taxas de transição}
%
Seja
%
\begin{equation}
  P_n(k)\equiv \Prob(N_1(k)=n),\qquad n=0,1,\ldots,N.
\end{equation}
%
A probabilidade de estar no estado $n$ após uma atualização recebe contribuições de apenas dois estados anteriores, $n-1$ e $n+1$, pois em cada passo exatamente uma bola troca de urna. A probabilidade de transição de um estado de origem $n$ para $n+1$ é
%
\begin{equation}
  T_{n\to n+1}=\frac{N-n}{N},
  \label{eq:Tup}
\end{equation}
%
pois é necessário selecionar uma das $N-n$ bolas da urna 2. Analogamente,
%
\begin{equation}
  T_{n\to n-1}=\frac{n}{N},
  \label{eq:Tdown}
\end{equation}
%
pois é necessário selecionar uma das $n$ bolas da urna 1. As duas probabilidades somam um:
\[
T_{n\to n+1}+T_{n\to n-1}=\frac{N-n}{N}+\frac{n}{N}=\frac{N}{N}= 1.
\]
%
A equação de evolução é, então,
%
\begin{align}
  P_n(k+1)
  &=P_{n-1}(k)T_{n-1\to n}+P_{n+1}(k)T_{n+1\to n}\nonumber\\
  &=P_{n-1}(k)\frac{N-(n-1)}{N}+P_{n+1}(k)\frac{n+1}{N}.
\end{align}
%
Portanto,
%
\begin{equation}
  \boxed{P_n(k+1)=P_{n-1}(k)\frac{N-n+1}{N}+P_{n+1}(k)\frac{n+1}{N}.}
  \label{eq:master}
\end{equation}
Adotamos a convenção $P_{-1}(k)=P_{N+1}(k)=0$, de modo que a mesma expressão também vale nos extremos $n=0$ e $n=N$. Na notação do enunciado, os fatores que multiplicam $P_{n-1}$ e $P_{n+1}$ são
%
\begin{equation}
  \boxed{W_1(n)=\frac{N-n+1}{N},\qquad W_2(n)=\frac{n+1}{N}.}
  \label{eq:W12}
\end{equation}
%
É importante interpretar corretamente os índices: $W_1(n)$ é a probabilidade de transição do estado $n-1$ para o estado $n$, enquanto $W_2(n)$ é a probabilidade de transição do estado $n+1$ para o estado $n$. Por isso, elas são funções do estado de destino $n$ nesta notação.

A formulação assume que, a cada passo, uma bola é escolhida uniformemente entre as $N$ bolas e sua localização é invertida. O estado $N_1(k)$ contém informação suficiente para determinar as probabilidades da próxima atualização; não é necessário conhecer a história anterior. Essa é a propriedade de Markov.

Como verificação, a equação conserva a normalização. Somando \eqref{eq:master} sobre $n$ e fazendo uma mudança de índice,
%
\begin{align}
  \sum_{n=0}^{N}P_n(k+1)
  &=\sum_{j=0}^{N}P_j(k)\frac{N-j}{N}
   +\sum_{j=0}^{N}P_j(k)\frac{j}{N}\nonumber\\
  &=\sum_{j=0}^{N}P_j(k)\frac{N-j+j}{N}\nonumber\\
  &=\sum_{j=0}^{N}P_j(k)=1.
\end{align}
%
Esse teste ajuda a evitar um erro comum: os dois fatores $W_1(n)$ e $W_2(n)$ não devem ser somados mantendo o mesmo índice de destino como se descrevessem duas saídas do mesmo estado. Para obter as probabilidades de saída de um estado $n$, devem-se usar $T_{n\to n-1}=n/N$ e $T_{n\to n+1}=(N-n)/N$.
%
\section*{(c) Distribuição estacionária e momentos}
%
Assumindo uma distribuição binomial,
%
\begin{equation}
  P_N (N_1)=\frac{1}{2^N}\binom{N}{N_1}=\frac{1}{2^N}\frac{N!}{N_1!(N-N_1)!}.
  \label{eq:pi}
\end{equation}
%
Antes de verificar que essa distribuição é uma solução da equação de evolução temporal no equilíbrio ($t \rightarrow \infty$), observamos que ela está normalizada, pois o teorema binomial fornece
%
\begin{equation}
  \sum_{n=0}^{N}P_N(N_1)
  =\frac{1}{2^N}\sum_{N_1=0}^{N}\binom{N}{N_1}
  =\frac{(1+1)^N}{2^N}=1.
\end{equation}
%
Para verificar que $P_N$ é estacionária, substituímos a distribuição na equação \eqref{eq:master}. Para $1\leq N_1\leq N-1$, o primeiro termo do lado direito é
%
\begin{align}
  P(N_1-1, t)W_1 = P(N_1-1,t)\frac{N-n+1}{N}
  &=\frac{1}{2^N}\frac{N!}{(N_1-1)!(N-N_1+1)!}\frac{N-N_1+1}{N}\nonumber\\
  &=\frac{1}{2^N}\frac{N!}{N}\frac{1}{(N_1-1)!(N-N_1)!}\nonumber\\
  &=\frac{1}{2^N}\frac{(N-1)!}{(N_1-1)!(N-N_1)!}\nonumber\\
  &=\frac{1}{2^N}\frac{N!}{N_1!(N-N_1)!}\frac{N_1}{N}\\
  &=P_{N}(N_1)\frac{N_1}{N}.
  \label{eq:firststationaryterm}
\end{align}
%
Na penúltima passagem usamos $N!=N(N-1)!$ e $N_1!=N_1(N_1-1)!$. De maneira análoga, o segundo termo é
%
\begin{align}
  P(N_1+1,t) W_2 = P(N_1+1,t) \frac{N_1+1}{N}
  &=\frac{1}{2^N}\frac{N!}{(N_1+1)!(N-N_1-1)!}\frac{N_1+1}{N}\nonumber\\
  &=\frac{1}{2^N}\frac{(N-1)!}{N_1!(N-N_1-1)!}\nonumber\\
  &=\frac{1}{2^N}\frac{N!}{N_1!(N-N_1)!}\frac{N-N_1}{N}\nonumber\\
  &=P_N(N_1)\frac{N-N_1}{N}.
  \label{eq:secondstationaryterm}
\end{align}
%
Somando as duas contribuições,
%
\begin{align}
  P(N_1, t + \Delta t) = P(N_1-1,t)W_1+P(N_1+1, t)W_2
  &=P_N(N_1)\frac{N_1}{N}+P_N(N_1)\frac{N-N_1}{N}\nonumber\\
  &=P_N(N_1)\frac{N_1+N-N_1}{N}\nonumber\\
  &=P_N(N_1).
\end{align}
%
Logo, $P_N(N_1, t + \Delta t)=P_N(N_1)$ é uma solução estacionária, ou seja, se o sistema for preparado segundo a distribuição binomial, essa distribuição permanece inalterada após cada atualização. Nos extremos $N_1=0$ e $N_1=N$, a mesma conclusão segue diretamente da equação de evolução com $P_{-1}=P_{N+1}=0$.
%
Agora, para calcular expressões para os valores médios ou esperados, por definição,
%
\begin{align}
  \avg{N_1}
  &=\sum_{N_1=0}^{N}N_1P_n
   =\frac{1}{2^N}\sum_{N_1=0}^{N}N_1\binom{N}{N_1}.
\end{align}
%
Usamos a identidade combinatória
%
\begin{equation}
  N_1\binom{N}{N_1}=N\binom{N-1}{N_1-1},
  \label{eq:combmean}
\end{equation}
%
que pode ser verificada diretamente:
\[
 N_1\binom{N}{N_1}=N_1\frac{N!}{N_1!(N-N_1)!}
 =\frac{N!}{(N_1-1)!(N-N_1)!}
 =N\frac{(N-1)!}{(N_1-1)!(N-N_1)!}.
\]
Assim,
%
\begin{align}
  \avg{N_1}
  &=\frac{N}{2^N}\sum_{N_1=1}^{N}\binom{N-1}{N_1-1}.
\end{align}
%
Definindo $j=N_1-1$, a soma torna-se
%
\begin{equation}
  \sum_{N_1=1}^{N}\binom{N-1}{N_1-1}
  =\sum_{j=0}^{N-1}\binom{N-1}{j}
  =(1+1)^{N-1}=2^{N-1}.
\end{equation}
%
Portanto,
%
\begin{equation}
  \boxed{\avg{N_1}=\frac{N}{2^N}2^{N-1}=\frac{N}{2}.}
  \label{eq:mean}
\end{equation}
%
Antes de calcular o desvio quadrático médio ou variância, é conveniente primeiro calcular o momento fatorial $\avg{N_1(N_1-1)}$. Temos
%
\begin{equation}
  \avg{N_1(N_1-1)}
  =\frac{1}{2^N}\sum_{N_1=0}^{N}N_1(N_1-1)\binom{N}{N_1}.
\end{equation}
%
A identidade necessária é
%
\begin{equation}
  N_1(N_1-1)\binom{N}{N_1}=N(N-1)\binom{N-2}{N_1-2}.
  \label{eq:combsecond}
\end{equation}
%
De fato,
%
\begin{align*}
  N_1(N_1-1)\binom{N}{N_1}
  &=N_1(N_1-1)\frac{N!}{N_1!(N-N_1)!}
   =\frac{N!}{(N_1-2)!(N-N_1)!}\\
  &=N(N-1)\frac{(N-2)!}{(N_1-2)!(N-N_1)!}
   =N(N-1)\binom{N-2}{N_1-2}.
\end{align*}
%
Logo,
%
\begin{align}
  \avg{N_1(N_1-1)}
  &=\frac{N(N-1)}{2^N}\sum_{N_1=2}^{N}\binom{N-2}{N_1-2}\nonumber\\
  &=\frac{N(N-1)}{2^N}\sum_{j=0}^{N-2}\binom{N-2}{j}\nonumber\\
  &=\frac{N(N-1)}{2^N}2^{N-2}\nonumber\\
  &=\frac{N(N-1)}{4}.
  \label{eq:factorialmoment}
\end{align}
%
Como $N_1^2=N_1(N_1-1)+N_1$, segue que
%
\begin{align}
  \avg{N_1^2}
  &=\avg{N_1(N_1-1)}+\avg{N_1}\nonumber\\
  &=\frac{N(N-1)}{4}+\frac{N}{2}\nonumber\\
  &=\frac{N^2-N+2N}{4}\nonumber\\
  &=\frac{N^2+N}{4}.
  \label{eq:secondmoment}
\end{align}
%
A variância é
%
\begin{align}
  \Var(N_1) = \sigma^2
  &=\avg{(N_1 - \avg{N_1})^2}\nonumber\\
  &=\avg{N_1^2}-\avg{N_1}^2\nonumber\\
  &=\frac{N^2+N}{4}-\left(\frac{N}{2}\right)^2\nonumber\\
  &=\frac{N^2+N}{4}-\frac{N^2}{4}\nonumber\\
  &=\boxed{\frac{N}{4}}.
\end{align}
%
Portanto, $\sigma=\sqrt{N}/2$ e $\sigma/\avg{N_1}= \frac{\sqrt{N}/2}{N/2}= 1/\sqrt{N}$, como antecipado no item (a).
%
\section*{(d) Evolução temporal do valor esperado}
%
Definimos o valor esperado no passo $k$ por
%
\begin{equation}
  m_k\equiv\avg{N_1}_k=\sum_{n=0}^{N}nP_n(k).
  \label{eq:mk}
\end{equation}
Queremos obter uma equação de recorrência para $m_k$ a partir da equação mestra \eqref{eq:master}. Usando \eqref{eq:master}, temos
%
\begin{align}
  m_{k+1}
  &=\sum_{n=0}^{N}nP_n(k+1)\nonumber\\
  &=\sum_{n=0}^{N}nP_{n-1}(k)\frac{N-n+1}{N}
   +\sum_{n=0}^{N}nP_{n+1}(k)\frac{n+1}{N}.
  \label{eq:meanstart}
\end{align}
%
Vamos tratar as duas somas separadamente, cuidando dos índices. Na primeira soma, fazemos $j=n-1$, isto é, $n=j+1$. Como $n$ varia de $1$ a $N$ para que $P_{n-1}$ seja não nulo, temos
%
\begin{align}
  S_1
  &=\sum_{n=1}^{N}nP_{n-1}(k)\frac{N-n+1}{N}\nonumber\\
  &=\sum_{j=0}^{N-1}(j+1)P_j(k)\frac{N-(j+1)+1}{N}\nonumber\\
  &=\sum_{j=0}^{N-1}(j+1)P_j(k)\frac{N-j}{N}.
  \label{eq:S1}
\end{align}
%
Na segunda soma, fazemos $j=n+1$, isto é, $n=j-1$. Como $n$ varia de $0$ a $N-1$ para que $P_{n+1}$ seja não nulo,
%
\begin{align}
  S_2
  &=\sum_{n=0}^{N-1}nP_{n+1}(k)\frac{n+1}{N}\nonumber\\
  &=\sum_{j=1}^{N}(j-1)P_j(k)\frac{j}{N}.
  \label{eq:S2}
\end{align}
%
Combinando \eqref{eq:S1} e \eqref{eq:S2}, podemos escrever
%
\begin{align}
  m_{k+1}
  &=\frac{1}{N}\sum_{j=0}^{N}P_j(k)
       \left[(j+1)(N-j)+(j-1)j\right].
  \label{eq:meancombined}
\end{align}
%
Os termos que foram estendidos aos extremos $j=0$ e $j=N$ são nulos quando deveriam estar ausentes: em $j=0$, $(j-1)j=0$; em $j=N$, $(j+1)(N-j)=0$. Expandimos cada produto entre colchetes:
%
\begin{align*}
  (j+1)(N-j)&=jN-j^2+N-j,\\
  (j-1)j&=j^2-j.
\end{align*}
%
Somando,
%
\begin{align}
  (j+1)(N-j)+(j-1)j
  &=jN-j^2+N-j+j^2-j\nonumber\\
  &=jN+N-2j\nonumber\\
  &=N+j(N-2).
  \label{eq:bracketidentity}
\end{align}
%
Substituindo \eqref{eq:bracketidentity} em \eqref{eq:meancombined},
%
\begin{align}
  m_{k+1}
  &=\frac{1}{N}\sum_{j=0}^{N}P_j(k)\left[N+j(N-2)\right]\nonumber\\
  &=\frac{1}{N}\left[N\sum_{j=0}^{N}P_j(k)+(N-2)\sum_{j=0}^{N}jP_j(k)\right].
\end{align}
%
Usamos agora a normalização $\sum_j P_j(k)=1$ e a definição de média $\sum_j jP_j(k)=m_k$:
%
\begin{align}
  m_{k+1}
  &=\frac{1}{N}\left[N+(N-2)m_k\right]\nonumber\\
  &=1+\frac{N-2}{N}m_k.
\end{align}
%
Obtemos a recorrência fundamental:
%
\begin{equation}
  \boxed{m_{k+1}=1+\left(1-\frac{2}{N}\right)m_k.}
  \label{eq:meanrecurrence}
\end{equation}
%
No regime estacionário da média, $m_{k+1}=m_k=m_\infty$. Substituindo na recorrência,
%
\begin{align}
  m_\infty&=1+\left(1-\frac{2}{N}\right)m_\infty,\\
  m_\infty-\left(1-\frac{2}{N}\right)m_\infty&=1,\\
  \left[1-1+\frac{2}{N}\right]m_\infty&=1,\\
  \frac{2}{N}m_\infty&=1.
\end{align}
%
Consequentemente,
%
\begin{equation}
  {m_\infty=\frac{N}{2},}
\end{equation}
%
em concordância com o resultado do item (c). Para resolver a recorrência, definimos o desvio em relação ao valor fixo,
%
\begin{equation}
  \delta_k\equiv m_k-\frac{N}{2}.
\end{equation}
%
Então
%
\begin{align}
  \delta_{k+1}
  &=m_{k+1}-\frac{N}{2}\nonumber\\
  &=1+\left(1-\frac{2}{N}\right)m_k-\frac{N}{2}\nonumber\\
  &=\left(1-\frac{2}{N}\right)m_k+\left(1-\frac{N}{2}\right).
\end{align}
%
Como
\[
1-\frac{N}{2}=-\frac{N-2}{2}=-\left(1-\frac{2}{N}\right)\frac{N}{2},
\]
podemos fatorar:
%
\begin{align}
  \delta_{k+1}
  &=\left(1-\frac{2}{N}\right)\left(m_k-\frac{N}{2}\right)\nonumber\\
  &=\left(1-\frac{2}{N}\right)\delta_k.
\end{align}
%
Iterando essa relação,
%
\begin{equation}
  \delta_k=\left(1-\frac{2}{N}\right)^k\delta_0,
\end{equation}
%
onde $\delta_0=m_0-N/2$. Portanto, a solução geral é
%
\begin{equation}
  \boxed{\avg{N_1}_k=\frac{N}{2}+\left(m_0-\frac{N}{2}\right)\left(1-\frac{2}{N}\right)^k.}
  \label{eq:generalsolution}
\end{equation}
%
Se inicialmente todas as bolas estão na urna 1, $m_0=N$, de modo que
%
\begin{align}
  \avg{N_1}_k
  &=\frac{N}{2}+\left(N-\frac{N}{2}\right)\left(1-\frac{2}{N}\right)^k\nonumber\\
  &=\frac{N}{2}+\frac{N}{2}\left(1-\frac{2}{N}\right)^k.
\end{align}
%
Logo,
%
\begin{equation}
  \boxed{\avg{N_1}_k=\frac{N}{2}\left[1+\left(1-\frac{2}{N}\right)^k\right].}
  \label{eq:initialall}
\end{equation}
%
Para $N>2$, o fator $1-2/N$ está entre zero e um. Escrevemos sua potência como
\[
\left(1-\frac{2}{N}\right)^k=\exp\left[k\ln\left(1-\frac{2}{N}\right)\right].
\]
Definindo o tempo de relaxação em unidades de atualizações por
%
\begin{equation}
  \tau=\frac{-1}{\ln\left(1-\frac{2}{N}\right)},
  \label{eq:tau}
\end{equation}
%
podemos escrever o fator de decaimento como $\e{-k/\tau}$. Para $N\gg1$, a expansão $\ln(1-x)=-x+O(x^2)$ dá
%
\begin{equation}
  \ln\left(1-\frac{2}{N}\right)\simeq -\frac{2}{N},
  \qquad
  \tau\simeq\frac{N}{2}.
\end{equation}
%
Numericamente,
%
\begin{center}
\begin{tabular}{@{}ccc@{}}
\toprule
$N$ & $\tau=-1/\ln(1-2/N)$ & Aproximação $N/2$\\
\midrule
$20$ & $9.49\ldots$ & $10$\\
$200$ & $99.50\ldots$ & $100$\\
\bottomrule
\end{tabular}
\end{center}
%
Isso explica por que o tempo necessário para relaxar cresce proporcionalmente a $N$. Para comparar sistemas com diferentes tamanhos, representar $N_1/N$ em função de $k/N$ é natural. As flutuações relativas são menores para $N=200$, embora as flutuações absolutas sejam maiores.
% ============================================================
% PROBLEMA 2
% ============================================================
\section*{Problema 2 -- Algoritmo de Metropolis para o modelo de Ising}
%
Consideramos o modelo de Ising ferromagnético, a campo externo nulo, em uma rede quadrada de tamanho linear $N$ e, portanto, com $N_s=N^2$ sítios. Em cada sítio há uma variável de Ising
\[
    \sigma_i\in\{-1,+1\},\qquad J>0,
\]
e adotaremos condições periódicas de contorno, de modo que todo sítio tenha quatro vizinhos mais próximos. A temperatura reduzida será denotada por
\[
    \thetaT\equiv \frac{k_B T}{J},
\]
para não confundir a temperatura física $T$ com a temperatura expressa em unidades de $J/k_B$. Também escreveremos
\[
    \beta=\frac{1}{k_BT},\qquad \beta J=\frac{1}{\thetaT}.
\]
Agora, vamos implementar o algoritmo de Metropolis ara o modelo de Ising de spin-$1/2$ a campo nulo, numa rede quadrada $N \times N$ , a temperatura $T$. O Hamiltoniano, sem campo externo, é
%
\begin{equation}
    \mathcal H(\bm\sigma)
    =-J\sum_{\langle ij\rangle}\sigma_i\sigma_j,
    \label{eq:hamiltoniano}
\end{equation}
%
onde $\langle ij\rangle$ indica que a soma é feita uma única vez sobre cada par de vizinhos mais próximos. Suponha que o algoritmo selecione o sítio $i$ e proponha inverter o valor do spin neste sítio,
%
\begin{equation}
    \sigma_i\longrightarrow \sigma_i'=-\sigma_i,
\end{equation}
%
mantendo todos os demais spins inalterados. Queremos obter
%
\[
    \Delta E\equiv \mathcal H(\bm\sigma')-\mathcal H(\bm\sigma).
\]
%
Como somente o spin $i$ muda, todos os termos do Hamiltoniano que não contêm $\sigma_i$ cancelam-se na diferença. A parte da energia associada às quatro ligações que chegam ao sítio $i$, antes da inversão, é
%
\begin{equation}
    E_i^{\mathrm{antes}}
    =-J\sigma_i\sum_{j\in\nn(i)}\sigma_j.
    \label{eq:Eantes}
\end{equation}
%
Depois da inversão, substituímos $\sigma_i$ por $-\sigma_i$:
%
\begin{align}
    E_i^{\mathrm{depois}}
    &=-J(-\sigma_i)\sum_{j\in\nn(i)}\sigma_j\\
    &=+J\sigma_i\sum_{j\in\nn(i)}\sigma_j.
    \label{eq:Edepois}
\end{align}
%
Subtraindo a energia inicial da energia final,
%
\begin{align}
    \Delta E
    &=E_i^{\mathrm{depois}}-E_i^{\mathrm{antes}}\\
    &=J\sigma_i\sum_{j\in\nn(i)}\sigma_j
      -\left(-J\sigma_i\sum_{j\in\nn(i)}\sigma_j\right)\\
    &=J\sigma_i\sum_{j\in\nn(i)}\sigma_j
      +J\sigma_i\sum_{j\in\nn(i)}\sigma_j\\
    &=2J\sigma_i\sum_{j\in\nn(i)}\sigma_j.
\end{align}
%
Portanto,
%
\begin{equation}
    \Delta E=2J\sigma_i\sum_{j\in\nn(i)}\sigma_j.
    \label{eq:deltaE}
\end{equation}
%
Para abreviar, definimos a soma dos quatro spins vizinhos
%
\begin{equation}
    S_i\equiv\sum_{j\in\nn(i)}\sigma_j.
    \label{eq:Si}
\end{equation}
%
Como cada spin vizinho vale $+1$ ou $-1$, se $n_+$ dos quatro vizinhos estiverem para cima, então haverá $4-n_+$ spins para baixo e
%
\begin{align}
    S_i&=(+1)n_+ +(-1)(4-n_+)\\
       &=n_+-4+n_+\\
       &=2n_+-4.
\end{align}
%
Como $n_+=0,1,2,3,4$, segue que
%
\begin{equation}
    S_i\in\{-4,-2,0,2,4\}.
\end{equation}
%
A equação \eqref{eq:deltaE} pode então ser escrita como
%
\begin{equation}
    \frac{\Delta E}{J}=2\sigma_i S_i.
\end{equation}
%
Uma vez que $\sigma_i=\pm1$, os valores possíveis são
%
\begin{equation}
    \frac{\Delta E}{J}\in\{-8,-4,0,4,8\}.
    \label{eq:deltaEpossiveis}
\end{equation}
%
A tabela abaixo mostra explicitamente como os valores de $S_i$ e do spin central determinam a variação de energia. A última coluna já contém a probabilidade de aceitação que será deduzida na próxima subseção.
%
\begin{table}[H]
\centering
\caption{Variação de energia para uma inversão local. Usamos $\thetaT=k_BT/J$.}
\label{tab:deltaE}
\begin{tabular}{rrrrl}
\toprule
$S_i$ & $\sigma_i$ & $\sigma_iS_i$ & $\Delta E/J$ & $p_{\mathrm{aceitar}}$\\
\midrule
$+4$ & $+1$ & $+4$ & $+8$ & $\e{{-8/\thetaT}}$\\
$+4$ & $-1$ & $-4$ & $-8$ & $1$\\
$+2$ & $+1$ & $+2$ & $+4$ & $\e{{-4/\thetaT}}$\\
$+2$ & $-1$ & $-2$ & $-4$ & $1$\\
$0$  & $\pm1$& $0$  & $0$  & $1$\\
$-2$ & $+1$ & $-2$ & $-4$ & $1$\\
$-2$ & $-1$ & $+2$ & $+4$ & $\e{-4/\thetaT}$\\
$-4$ & $+1$ & $-4$ & $-8$ & $1$\\
$-4$ & $-1$ & $+4$ & $+8$ & $\e{-8/\thetaT}$\\
\bottomrule
\end{tabular}
\end{table}
%
Observe a interpretação física:
%
\begin{itemize}
    \item se $\Delta E<0$, a inversão diminui a energia e será aceita;
    \item se $\Delta E=0$, a energia não muda e a inversão também será aceita;
    \item se $\Delta E>0$, a inversão aumenta a energia e só será aceita com uma probabilidade menor que um.
\end{itemize}
%
A regra de Metropolis é
%
\begin{equation}
    p_{\mathrm{aceitar}}(\bm\sigma\to\bm\sigma')
    =\min\left\{1,\e{-\beta\Delta E}\right\}.
    \label{eq:metropolisbeta}
\end{equation}
%
Isso equivale à forma por partes
%
\begin{equation}
    p_{\mathrm{aceitar}}=
    \begin{cases}
        1, & \Delta E\leq0,\\[2mm]
        \e{-\beta\Delta E}, & \Delta E>0.
    \end{cases}
\end{equation}
%
Para escrever a expressão em unidades de $J$, definimos
%
\[
    \epsilon\equiv\frac{\Delta E}{J},\qquad
    \thetaT\equiv\frac{k_BT}{J}.
\]
%
Então
%
\begin{align}
    \beta\Delta E
    &=\frac{1}{k_BT}\Delta E\\
    &=\frac{\Delta E/J}{k_BT/J}\\
    &=\frac{\epsilon}{\thetaT}.
\end{align}
%
Logo, a expressão que pode ser usada diretamente no programa é
%
\begin{equation}
    p_{\mathrm{aceitar}}=\min\left\{1,\e{-\epsilon/\thetaT}\right\}.
    \label{eq:metropolis_theta}
\end{equation}
%
Em particular, como $\epsilon$ só pode assumir os valores $-8,-4,0,4,8$, só é necessário calcular os dois números
%
\begin{equation}
    p_4=\e{-4/\thetaT},\qquad p_8=\e{-8/\thetaT}
\end{equation}
%
para cada temperatura. Os casos com energia não crescente são aceitos com probabilidade um. Para $\Delta E>0$, sorteia-se $z$ uniformemente em $[0,1)$. A inversão é aceita se
%
\begin{equation}
    z<p_{\mathrm{aceitar}}=\e{-\Delta E/(k_BT)}.
\end{equation}
%
A probabilidade de que isso aconteça é precisamente $p_{\mathrm{aceitar}}$, pois $z$ é uniforme no intervalo unitário. Se $z$ for maior ou igual a essa probabilidade, conserva-se o spin original. A distribuição de equilíbrio que queremos amostrar é
%
\begin{equation}
    P(\bm\sigma)=\frac{1}{Z}\e{-\beta\mathcal H(\bm\sigma)},
    \qquad
    Z=\sum_{\{\bm\sigma\}}\e{-\beta\mathcal H(\bm\sigma)}.
    \label{eq:boltzmann}
\end{equation}
%
Aqui, $Z$ é a função de partição e a soma percorre todas as $2^{N^2}$ configurações de spins possíveis. Considere duas configurações $\bm\sigma$ e $\bm\sigma'$ que diferem apenas pela inversão de um spin. A variação de energia é
%
\[
    \Delta E=\mathcal H(\bm\sigma')-\mathcal H(\bm\sigma).
\]
%
Se a proposta de inverter o spin é simétrica, isto é, a probabilidade de propor a transição é igual à de propor a transição reversa, a regra de Metropolis satisfaz o balanço detalhado. Vamos verificar isso sem omitir casos.

\noindent\textbf{Caso 1: $\Delta E>0$.} A transição direta aumenta a energia, portanto
%
\[
    A(\bm\sigma\to\bm\sigma')=\e{-\beta\Delta E}.
\]
%
Na transição inversa, a variação de energia é $-\Delta E<0$, logo a aceitação é um:
%
\[
    A(\bm\sigma'\to\bm\sigma)=1.
\]
%
Assim,
%
\begin{equation}
    \frac{A(\bm\sigma\to\bm\sigma')}{A(\bm\sigma'\to\bm\sigma)}
    =\e{-\beta\Delta E}.
\end{equation}
%
\noindent\textbf{Caso 2: $\Delta E<0$.} A transição direta diminui a energia e é aceita com probabilidade um:
%
\[
    A(\bm\sigma\to\bm\sigma')=1.
\]
%
A transição inversa aumenta a energia em $-\Delta E>0$, então
%
\[
    A(\bm\sigma'\to\bm\sigma)=\e{-\beta(-\Delta E)}=\e{\beta\Delta E}.
\]
%
Consequentemente,
%
\begin{equation}
    \frac{A(\bm\sigma\to\bm\sigma')}{A(\bm\sigma'\to\bm\sigma)}
    =\frac{1}{\e{\beta\Delta E}}=\e{-\beta\Delta E}.
\end{equation}
%
\noindent\textbf{Caso 3: $\Delta E=0$.} As duas transições são aceitas com probabilidade um; a razão é um, igual a $\e0$. Portanto, nos três casos,
%
\begin{equation}
    \frac{A(\bm\sigma\to\bm\sigma')}{A(\bm\sigma'\to\bm\sigma)}
    =\e{-\beta\Delta E}.
    \label{eq:razao_aceitacao}
\end{equation}
%
Agora calculemos a razão dos pesos de Boltzmann:
%
\begin{align}
    \frac{P(\bm\sigma')}{P(\bm\sigma)}
    &=\frac{Z^{-1}\e{-\beta\mathcal H(\bm\sigma')}}
            {Z^{-1}\e{-\beta\mathcal H(\bm\sigma)}}\\
    &=\e{-\beta[\mathcal H(\bm\sigma')-\mathcal H(\bm\sigma)]}\\
    &=\e{-\beta\Delta E}.
    \label{eq:razao_boltzmann}
\end{align}
%
Comparando \eqref{eq:razao_aceitacao} e \eqref{eq:razao_boltzmann}, obtemos
%
\[
    \frac{P(\bm\sigma')}{P(\bm\sigma)}
    =\frac{A(\bm\sigma\to\bm\sigma')}{A(\bm\sigma'\to\bm\sigma)}.
\]
%
Multiplicando cruzado,
%
\begin{equation}
    P(\bm\sigma)A(\bm\sigma\to\bm\sigma')
    =P(\bm\sigma')A(\bm\sigma'\to\bm\sigma).
    \label{eq:balance_detalhado}
\end{equation}
%
Esta é a condição de balanço detalhado para a atualização local (o fator da proposta simétrica cancela dos dois lados). Ela implica que a distribuição de Boltzmann é estacionária para cada atualização local. Uma varredura sistemática é uma composição dessas atualizações; cada atualização preserva a distribuição de Boltzmann, e por isso uma varredura completa também a preserva, ainda que a composição ordenada não precise ser reversível como uma única transição global. Juntamente com a ergodicidade (toda configuração é alcançável a partir de qualquer outra por uma sequência finita de inversões), isso garante que a distribuição de Boltzmann é a
\emph{única} distribuição estacionária da cadeia, e que a cadeia converge para ela independentemente da configuração inicial~\cite{NewmanBarkema1999,LandauBinder2005}. Uma implementação que segue a ordem sequencial descrita no enunciado é:
%
\begin{enumerate}
    \item Inicializar a matriz de spins, por exemplo, com $\sigma_{ij}=+1$ em todos os sítios.
    \item Percorrer os sítios em uma ordem fixa. Um \emph{passo de Monte Carlo} (MCS, do inglês \textit{Monte Carlo step}) consiste em visitar todos os $N^2$ sítios uma vez.
    \item Para cada sítio $(i,j)$, somar os quatro vizinhos e calcular
    \[
        \epsilon=\frac{\Delta E}{J}=2\sigma_{ij}
        (\sigma_{i+1,j}+\sigma_{i-1,j}+\sigma_{i,j+1}+\sigma_{i,j-1}),
    \]
    com os índices tomados periodicamente.
    \item Se $\epsilon\leq0$, inverter o spin. Se $\epsilon>0$, inverter com probabilidade $\exp(-\epsilon/\thetaT)$, usando um número aleatório uniforme.
    \item Repetir muitas varreduras. As primeiras são usadas para termalização e descartadas; as seguintes são usadas para medir observáveis.
\end{enumerate}
%
O enunciado pede uma ordem sequencial. Em outras implementações, escolhem-se sítios aleatoriamente. Nos dois casos, a regra local de aceitação deve ser implementada com a mesma variação de energia; a escolha concreta do esquema de atualização afeta as correlações temporais e a eficiência da simulação.
% ============================================================
Após a termalização, a cadeia de Metropolis gera configurações distribuídas de acordo com a medida de Boltzmann \eqref{eq:boltzmann}. Não é necessário que uma configuração microscópica específica seja visitada muitas vezes: em sistemas grandes, até configurações típicas individualmente possuem probabilidades muito pequenas. O sentido de ``representativas'' é que as configurações amostradas reproduzem as propriedades estatísticas de equilíbrio, como energia média, magnetização e suas flutuações.
%
Os regimes físicos são:
%
\begin{itemize}
    \item \textbf{Baixa temperatura, $T<T_c$:} o acoplamento ferromagnético favorece alinhamento. As configurações típicas têm grandes regiões de spins iguais e, longe de $T_c$, o módulo da magnetização é elevado. Ainda há defeitos térmicos e paredes de domínio a temperatura não nula.
    \item \textbf{Alta temperatura, $T>T_c$:} a energia térmica compete fortemente com o alinhamento. As configurações são desordenadas e a magnetização por sítio tende a zero quando o tamanho do sistema cresce.
    \item \textbf{Próximo de $T_c$:} surgem flutuações de grande escala e domínios correlacionados. O sistema relaxa lentamente e são necessárias séries maiores para obter médias confiáveis.
\end{itemize}
%
A magnetização por sítio de uma configuração é
%
\begin{equation}
    m(\bm\sigma)=\frac{1}{N^2}\sum_{i=1}^{N^2}\sigma_i.
    \label{eq:magnetizacao}
\end{equation}
%
A transformação global $\sigma_i\mapsto-\sigma_i$ deixa o Hamiltoniano sem campo inalterado, pois cada produto de vizinhos satisfaz
%
\[
    (-\sigma_i)(-\sigma_j)=\sigma_i\sigma_j.
\]
%
Logo, cada configuração de magnetização $m$ tem uma configuração de mesma energia e magnetização $-m$. Como seus pesos de Boltzmann são iguais, para uma rede finita a campo nulo resulta exatamente
%
\begin{equation}
    \langle m\rangle=0.
\end{equation}
%
Isso vale inclusive abaixo de $T_c$ se o ensemble de equilíbrio incluir ambos os setores, com magnetizações positiva e negativa. Para caracterizar a ordem ferromagnética em uma simulação finita, mede-se portanto o módulo:
%
\begin{equation}
    \langle |m|\rangle\simeq\frac{1}{M}\sum_{k=1}^{M}|m_k|,
    \label{eq:media_modulo}
\end{equation}
%
onde $m_k$ é a magnetização da $k$-ésima configuração amostrada após a termalização e $M$ é o número de medidas. O \emph{número de medidas} $M$ é o número de varreduras realizadas \emph{após} a termalização, durante as quais se coletam os valores dos observáveis
(magnetização, energia, etc.) a cada varredura. É importante distinguir duas afirmações: para qualquer $N$ finito e campo nulo, a média simétrica $\langle m\rangle$ é zero; já o limite termodinâmico pode apresentar magnetização espontânea não nula para $T<T_c$, quando uma fase é selecionada. O uso de $|m|$ é uma maneira prática de acompanhar a ordem sem cancelar contribuições de sinais opostos. Uma simulação atinge um regime compatível com equilíbrio quando seus observáveis já não mostram uma deriva sistemática causada pela condição inicial e passam a flutuar em torno de médias estáveis. Na prática, monitoram-se ao menos
%
\[
    m_k,\qquad |m_k|,\qquad e_k\equiv\frac{\mathcal H_k}{N^2},
\]
%
ao longo das varreduras.
%
A estabilização visual não é, sozinha, prova suficiente de equilíbrio. A estabilização visual não é, sozinha, prova suficiente de equilíbrio. Recomenda-se o
seguinte protocolo~\cite{NewmanBarkema1999,LandauBinder2005}:
%
\begin{enumerate}
    \item Descartar um intervalo inicial de termalização de $N_{\mathrm{term}}$ MCS.
    \item Comparar séries que partem de condições iniciais distintas (todos os spins para
          cima e spins aleatórios); se convergirem para as mesmas médias, há evidência de
          equilíbrio.
    \item Repetir a simulação com diferentes sementes aleatórias e verificar a
          reprodutibilidade dos resultados.
    \item Verificar a estabilidade das médias ao dobrar sucessivamente $N_{\mathrm{term}}$
          e $M$: se as médias não se alteram significativamente, a simulação é
          suficientemente longa.
    \item Estimar a autocorrelação temporal $\tau$ e os erros estatísticos por blocagem,
          Jackknife ou Bootstrap.
\end{enumerate}
%
Configurações consecutivas não são independentes. Se a magnetização muda lentamente, contar cada varredura como uma medida independente subestima o erro. O número efetivo de amostras independentes pode ser muito menor que $M$. Perto da temperatura crítica, o tempo de relaxação cresce com o tamanho do sistema. Para dinâmica local de Metropolis no modelo de Ising bidimensional, o comportamento assintótico é frequentemente descrito por
%
\begin{equation}
    \tau\sim |t|^{-z\nu},\qquad t=\frac{T-T_c}{T_c},
\end{equation}
%
com $\nu=1$ e um expoente dinâmico $z$ aproximadamente igual a $2.17$ para esta classe de atualização local. Esse expoente é dinâmico e depende do algoritmo; não deve ser confundido com o expoente estático $\nu$.
%
Para o modelo de Ising bidimensional na rede quadrada, no limite termodinâmico e a campo
nulo, a temperatura crítica exata foi obtida por Onsager em 1944~\cite{Onsager1944}
(com a relação de dualidade de Kramers--Wannier~\cite{KramersWannier1941}):
%
\begin{equation}
    \frac{k_BT_c}{J}=\thetaT_c
    =\frac{2}{\ln(1+\sqrt{2})}\simeq 2.269185.
    \label{eq:Tc}
\end{equation}
%
Em uma rede finita, não há uma singularidade termodinâmica verdadeira, as grandezas variam suavemente com a temperatura e a região crítica é arredondada. Para uma rede de tamanho linear $N$, a escala característica de arredondamento obedece, em termos de escala,
%
\begin{equation}
    \Delta T\sim N^{-1/\nu};
\end{equation}
%
para o Ising 2D, $\nu=1$, então essa escala diminui como $N^{-1}$. Ao aumentar $N$, o comportamento crítico se torna mais nítido, mas a magnetização não sofre um salto finito no modelo bidimensional, sua magnetização espontânea tende continuamente a zero quando $T\to T_c^-$. Como referência teórica para conferir a simulação, a magnetização espontânea exata no
limite termodinâmico foi derivada por Yang em 1952~\cite{Yang1952}:
%
\begin{equation}
    m_{\infty}(\thetaT)=
    \begin{cases}
    \left[1-\sinh^{-4}\!\left(\dfrac{2}{\thetaT}\right)\right]^{1/8},
        & \thetaT<\thetaT_c,\\[3mm]
    0, & \thetaT\geq\thetaT_c.
    \end{cases}
    \label{eq:magnetizacao_exata}
\end{equation}
%
Para tamanho finito, $\langle|m|\rangle$ geralmente permanece positivo mesmo acima de $T_c$; por isso os pontos simulados não precisam coincidir com a curva do limite termodinâmico. Para cada temperatura reduzida $\thetaT$:
%
\begin{enumerate}
    \item inicialize a rede;
    \item execute $N_{\mathrm{term}}$ MCS de termalização e descarte as configurações produzidas nesse intervalo;
    \item execute mais $M$ MCS e calcule $m_k$ pela equação \eqref{eq:magnetizacao};
    \item calcule a média de $|m|$ pela equação \eqref{eq:media_modulo};
    \item repita para uma faixa de temperaturas que cubra valores abaixo e acima de $\thetaT_c$;
    \item represente $\langle |m|\rangle$ em função de $\thetaT$, idealmente com barras de erro estimadas por blocagem ou por réplicas independentes.
\end{enumerate}
%
A Figura~\ref{fig:magnetizacao} mostra um exemplo ilustrativo para uma rede de tamanho finito, junto da curva exata \eqref{eq:magnetizacao_exata}. Os parâmetros exatos da simulação e a semente aleatória estão documentados no script fornecido. As diferenças perto de $T_c$ são esperadas para uma rede finita e para uma simulação de duração limitada; antes de usar dados como estimativas quantitativas, é necessário testar a convergência e aumentar os tempos de termalização e medição.
%
\begin{figure}[H]
    \centering
    \includegraphics[width=0.92\linewidth]{Figures/statistical_mechanics/monte_carlo.png}
    \caption{$\langle|m|\rangle$ obtido por Metropolis para redes $N\times N$ de tamanhos
    $N=4,10,20,30,40,60,100$, comparado à magnetização espontânea exata no limite
    termodinâmico (curva de Yang~\cite{Yang1952}). O eixo horizontal é $\widetilde T=
    k_BT/J$, a temperatura reduzida, equivalente à notação $kT/K_{\mathrm{nn}}$ de
    Landau~\cite{Landau1976}. As curvas simuladas acima de $\widetilde T_c$ não são nulas
    porque o tamanho é finito e a média é feita do módulo. As barras de erro são estimadas
    por blocagem; na região crítica, autocorrelações longas exigem testes adicionais de
    convergência~\cite{NightingaleBlote1996}. O código que gera esta figura está disponível
    no repositório do \href{https://github.com/ThiagoSDomingues/Statistical_Mechanics/tree/main/Listas/Lista_6}{Github}.}
    \label{fig:magnetizacao}
\end{figure}
%
O algoritmo de Metropolis descrito aqui neste problema foi implementado e o arquivo \texttt{metropolis\_ising.py} está disponível na minha página do \href{https://github.com/ThiagoSDomingues/Statistical_Mechanics/tree/main/Listas/Lista_6}{Github}. A Figura~\ref{fig:magnetizacao} reproduz o estilo do gráfico de Landau~\cite{Landau1976}
para redes de tamanhos $N=4,10,20,30,40,60,100$, comparando $\langle|m|\rangle$ com a
curva exata de Yang~\eqref{eq:magnetizacao_exata}. Este arquivo implementa a atualização sequencial com condições periódicas, descarta um intervalo de termalização e estima $\langle|m|\rangle$ e sua incerteza por blocos.  Aumentar o tamanho e o número de varreduras melhora alguns aspectos da análise, mas também aumenta o custo computacional. Em particular, o número de MCS necessário não é universal, deve ser determinado por testes de convergência para cada tamanho e faixa de temperatura. 
% ============================================================
% Referências
% ============================================================
\begin{thebibliography}{99}
 
\bibitem{Onsager1944}
L.~Onsager,
\textit{Crystal Statistics. I. A Two-Dimensional Model with an Order-Disorder Transition},
Physical Review \textbf{65}, 117 (1944).
\url{https://doi.org/10.1103/PhysRev.65.117}
 
\bibitem{KramersWannier1941}
H.~A.~Kramers e G.~H.~Wannier,
\textit{Statistics of the Two-Dimensional Ferromagnet. Part I},
Physical Review \textbf{60}, 252 (1941).
\url{https://doi.org/10.1103/PhysRev.60.252}
 
\bibitem{Yang1952}
C.~N.~Yang,
\textit{The Spontaneous Magnetization of a Two-Dimensional Ising Model},
Physical Review \textbf{85}, 808 (1952).
\url{https://doi.org/10.1103/PhysRev.85.808}
 
\bibitem{NightingaleBlote1996}
M.~P.~Nightingale e H.~W.~J.~Blöte,
\textit{Dynamic Exponent of the Two-Dimensional Ising Model and Monte Carlo Computation
of the Subdominant Eigenvalue of the Stochastic Matrix},
Physical Review Letters \textbf{76}, 4548 (1996).
\url{https://doi.org/10.1103/PhysRevLett.76.4548}
 
\bibitem{Landau1976}
D.~P.~Landau,
\textit{Finite-size behavior of the Ising square lattice},
Physical Review B \textbf{13}, 2997 (1976).
\url{https://doi.org/10.1103/PhysRevB.13.2997}
 
\bibitem{NewmanBarkema1999}
M.~E.~J.~Newman e G.~T.~Barkema,
\textit{Monte Carlo Methods in Statistical Physics},
Oxford University Press, Oxford (1999).
 
\bibitem{LandauBinder2005}
D.~P.~Landau e K.~Binder,
\textit{A Guide to Monte Carlo Simulations in Statistical Physics},
3ª~ed., Cambridge University Press (2009).

 
\end{thebibliography}

\end{document}
